\documentclass[UTF8,zihao=-4,fontset=windows]{ctexart}
\usepackage{amsmath,amssymb}
\usepackage[a4paper,margin=2.05cm]{geometry}
\usepackage{enumitem}
\setlist[enumerate]{itemsep=0.52em,topsep=0.28em,leftmargin=2.2em}
\setlength{\parindent}{0pt}
\setlength{\parskip}{0.4em}

\newcommand{\Prob}{\mathbb P}
\newcommand{\EX}{\mathbb E}
\newcommand{\DD}{\operatorname D}
\newcommand{\Var}{\operatorname{Var}}
\newcommand{\Cov}{\operatorname{Cov}}
\newcommand{\iid}{\overset{\mathrm{i.i.d.}}{\sim}}

\title{\vspace{-1.2cm}随机数学与统计期末仿真卷\\{\large 答案版}}
\date{}

\begin{document}
\maketitle
\vspace{-1.6em}

\section*{一、}
设 $m\ge 2$ 为正整数，$X,Y \iid U(0,m)$。
\begin{enumerate}
  \item 在正方形 $[0,m]^2$ 中，区域 $\{x+y\le 1\}$ 是直角三角形，面积为 $1/2$，故
  \[
    \Prob\{X+Y\le 1\}=\frac{1}{2m^2}.
  \]

  \item
  \[
    \Prob\{\min(X,Y)\le 1\}
    =1-\Prob\{X>1,Y>1\}
    =1-\left(\frac{m-1}{m}\right)^2
    =\frac{2m-1}{m^2}.
  \]

  \item 因 $m$ 为正整数，对 $0\le w\le 1$，
  \[
    \Prob\{W\le w\}
    =\sum_{k=0}^{m-1}\Prob\{k\le X<k+w\}
    =m\cdot\frac{w}{m}=w.
  \]
  故 $W\sim U(0,1)$。
\end{enumerate}

\section*{二、}
设 $X\sim \operatorname{Exp}(\lambda)$，且 $F_X(2\ln2)=3/4$。
\begin{enumerate}
  \item 由
  \[
    1-e^{-2\lambda\ln2}=\frac34
  \]
  得 $e^{-2\lambda\ln2}=1/4=2^{-2}$，故 $\lambda=1$。

  \item 此时 $\EX X=1,\ \DD X=1$，所以
  \[
    \Prob\{X-\EX X>\sqrt{\DD X}\}
    =\Prob\{X>2\}=e^{-2}.
  \]

  \item 记
  \[
    A_n=\frac{2S_n}{n(n+1)}.
  \]
  因 $\EX X_k=1$，
  \[
    \EX A_n=\frac{2}{n(n+1)}\sum_{k=1}^n k=1.
  \]
  又 $\Var(X_k)=1$ 且相互独立，
  \[
    \Var(A_n)=\frac{4}{n^2(n+1)^2}\sum_{k=1}^n k^2
    =\frac{4}{n^2(n+1)^2}\cdot\frac{n(n+1)(2n+1)}{6}\to 0.
  \]
  由 Chebyshev 不等式，$A_n\xrightarrow{P}1$。
\end{enumerate}

\section*{三、}
设 $\{B_t:t\ge0\}$ 为标准布朗运动。
\begin{enumerate}
  \item $U=2B_1+B_2-B_3$ 是正态随机变量，均值为 $0$。利用 $\Cov(B_s,B_t)=\min(s,t)$，
  \[
    \Var(U)=4+2+3+4-4-4=5.
  \]
  所以 $U\sim N(0,5)$，并且
  \[
    \Prob\{|U|\le\sqrt{\DD U}\}
    =\Prob\{|N(0,1)|\le 1\}
    =2\Phi(1)-1.
  \]

  \item 由独立增量，
  \[
    B_3\mid B_1=1\sim N(1,2),
  \]
  因而
  \[
    \Prob\{B_3<3\mid B_1=1\}=\Phi(\sqrt2).
  \]
  又 $(B_1,B_3)$ 为二维正态，
  \[
    B_1\mid B_3=3\sim N\left(\frac{\Cov(B_1,B_3)}{\Var(B_3)}3,\ 
    1-\frac{\Cov(B_1,B_3)^2}{\Var(B_3)}\right)
    =N\left(1,\frac23\right).
  \]
  故
  \[
    \Prob\{B_1<1\mid B_3=3\}=\frac12.
  \]
\end{enumerate}

\section*{四、}
联合密度为 $f_{X,Y}(x,y)=8xy,\ 0<y<x<1$。
\begin{enumerate}
  \item
  \[
    f_X(x)=\int_0^x 8xy\,dy=4x^3,\qquad 0<x<1.
  \]
  因此
  \[
    f_{Y\mid X=x}(y)=\frac{8xy}{4x^3}=\frac{2y}{x^2},\qquad 0<y<x.
  \]
  所以
  \[
    \EX(Y\mid X=x)=\int_0^x y\frac{2y}{x^2}\,dy=\frac{2x}{3},
    \qquad \EX(Y\mid X)=\frac{2X}{3}.
  \]

  \item 令 $Z=X-Y$。当 $0<z<1$ 时，
  \[
    f_Z(z)=\int_0^{1-z}8(z+y)y\,dy
    =4z(1-z)^2+\frac83(1-z)^3
    =\frac43(1-z)^2(z+2).
  \]
  因此
  \[
    f_Z(z)=
    \begin{cases}
      \dfrac43(1-z)^2(z+2), & 0<z<1,\\
      0, & \text{其他}.
    \end{cases}
  \]

  \item 记 $A=\{X+Y<1\}$。由于还需满足 $0<Y<X<1$，
  \[
    \Prob(A)=\int_0^{1/2}\int_0^x8xy\,dy\,dx+
    \int_{1/2}^1\int_0^{1-x}8xy\,dy\,dx=\frac16.
  \]
  同理
  \[
    \EX[X\mathbf 1_A]
    =\int_0^{1/2}\int_0^x8x^2y\,dy\,dx+
    \int_{1/2}^1\int_0^{1-x}8x^2y\,dy\,dx=\frac{11}{120}.
  \]
  故
  \[
    \EX(X\mid X+Y<1)=\frac{11/120}{1/6}=\frac{11}{20}.
  \]
\end{enumerate}

\section*{五、}
设 $X_1,\ldots,X_n \iid N(0,\sigma^2)$。
\begin{enumerate}
  \item 因 $\overline X\sim N(0,\sigma^2/n)$，
  \[
    M_{\overline X}(t)=\exp\left(\frac{\sigma^2t^2}{2n}\right),
    \qquad
    \varphi_{\overline X}(t)=\exp\left(-\frac{\sigma^2t^2}{2n}\right).
  \]

  \item 由正态样本的对称性或线性投影，
  \[
    \EX(X_i\mid \overline X)=\overline X,\qquad i=1,\ldots,n.
  \]
  因此
  \[
    \EX(2X_1+3X_2\mid \overline X)=5\overline X.
  \]

  \item $Y_1,Y_2,Y_3$ 均为零均值正态变量，且
  \[
    \Var(Y_1)=\Var(Y_2)=\Var(Y_3)=\sigma^2.
  \]
  直接计算可得 $\Cov(Y_i,Y_j)=0\ (i\ne j)$。联合正态下不相关即独立，故
  \[
    \frac{Y_1}{\sigma}\sim N(0,1),\qquad
    \frac{Y_2^2+Y_3^2}{\sigma^2}\sim \chi^2_2,
  \]
  且二者独立。因此
  \[
    \frac{Y_1}{\sqrt{(Y_2^2+Y_3^2)/2}}\sim t(2).
  \]

  \item 对 $X_i\sim N(0,\sigma^2)$，
  \[
    \EX|X_i|=\sigma\sqrt{\frac2\pi}.
  \]
  所以
  \[
    \EX T_n=\sqrt{\frac\pi2}\,\EX|X_i|=\sigma,
  \]
  即 $T_n$ 是 $\sigma$ 的无偏估计。
  又
  \[
    \Var(|X_i|)=\sigma^2\left(1-\frac2\pi\right).
  \]
  由中心极限定理，
  \[
    \sqrt n(T_n-\sigma)\xrightarrow{D}
    N\left(0,\frac\pi2\sigma^2\left(1-\frac2\pi\right)\right)
    =N\left(0,\sigma^2\left(\frac\pi2-1\right)\right).
  \]
\end{enumerate}

\section*{六、}
设 $f(x;\theta)=\theta x^{\theta-1},\ 0<x<1,\ \theta>0$，$T=-\sum_{i=1}^n\ln X_i$。
\begin{enumerate}
  \item 联合密度为
  \[
    \prod_{i=1}^n\theta x_i^{\theta-1}
    =\theta^n\exp\left((\theta-1)\sum_{i=1}^n\ln x_i\right)
    =\theta^n e^{-\theta T}\prod_{i=1}^n x_i^{-1}.
  \]
  由因子分解定理，$T$ 是 $\theta$ 的充分统计量。又 $T$ 的分布构成 Gamma 族
  \[
    f_T(t;\theta)=\frac{\theta^n}{\Gamma(n)}t^{n-1}e^{-\theta t},\qquad t>0,
  \]
  该一参数指数族对 $\theta>0$ 完全，故 $T$ 完全。

  \item 令 $Y_i=-\ln X_i$。则
  \[
    \Prob\{Y_i>y\}=\Prob\{X_i<e^{-y}\}=e^{-\theta y},\qquad y>0,
  \]
  所以 $Y_i\sim \operatorname{Exp}(\theta)$。因此
  \[
    T=\sum_{i=1}^nY_i\sim \Gamma(n,\text{rate }\theta),
  \]
  等价地，$2\theta T\sim \chi^2_{2n}$。

  \item 对数似然为
  \[
    \ell(\theta)=n\ln\theta-\theta T+\sum_{i=1}^n\ln x_i.
  \]
  求导得 $\ell'(\theta)=n/\theta-T$，故
  \[
    \hat\theta_{\mathrm{MLE}}=\frac{n}{T}.
  \]
  若 $\eta=1/\theta$，则由不变性
  \[
    \hat\eta_{\mathrm{MLE}}=\frac{T}{n}.
  \]
  且 $\EX(T/n)=1/\theta=\eta$。由于 $T$ 完全充分，$T/n$ 是 $\eta$ 的 UMVUE。

  \item 由 $2\theta T\sim \chi^2_{2n}$，
  \[
    \Prob\left\{\chi^2_{2n,\alpha/2}\le 2\theta T\le
    \chi^2_{2n,1-\alpha/2}\right\}=1-\alpha.
  \]
  所以 $\theta$ 的 $1-\alpha$ 双侧等尾置信区间为
  \[
    \left[
      \frac{\chi^2_{2n,\alpha/2}}{2T},
      \frac{\chi^2_{2n,1-\alpha/2}}{2T}
    \right].
  \]

  \item 对 $H_0:\theta=1$ 与 $H_1:\theta=2$，似然比
  \[
    \frac{L(2)}{L(1)}=2^n e^{-T}.
  \]
  该比值随 $T$ 单调递减，因此拒绝域形如 $\{T\le c_\alpha\}$。在 $H_0$ 下，
  $2T\sim\chi^2_{2n}$，故取
  \[
    c_\alpha=\frac12\chi^2_{2n,\alpha}.
  \]
  水平 $\alpha$ 的拒绝域为
  \[
    W=\left\{T\le \frac12\chi^2_{2n,\alpha}\right\}.
  \]
\end{enumerate}

\end{document}
